\section*{Input-output price shocks and inflation inequality}
\begin{table}[H]
\centering
\caption{Cumulative Energy Shock Exposure, June 2021--June 2023}
\scriptsize
\setlength{\tabcolsep}{1.2pt}
\begin{tabular}{lrrrrrrr}
\toprule
Country & \shortstack{Effect on average\\inflation -- Total} & \shortstack{Direct effect on\\average inflation} & \shortstack{Indirect effect on\\average inflation} & \shortstack{Indirect effect on average\\inflation through trade} & \shortstack{Effect on\\inflation gap} & \shortstack{Direct effect on\\inflation gap} & \shortstack{Indirect effect on\\inflation gap}
 \\
\midrule
Austria & 24.1 & 6.6 & 17.5 & 1.3 & 11.0 & 4.7 & 6.4
 \\
Belgium & 2.2 & 0.3 & 1.9 & 1.8 & -0.3 & -0.0 & -0.2
 \\
Cyprus & 4.2 & 2.1 & 2.0 & 0.4 & 2.4 & 2.2 & 0.2
 \\
Germany & 7.1 & 3.9 & 3.2 & 0.5 & 1.3 & 1.6 & -0.4
 \\
Estonia & 22.8 & 12.9 & 9.9 & 1.0 & 27.1 & 22.8 & 4.3
 \\
Greece & 2.9 & 1.5 & 1.4 & 0.3 & 1.4 & 1.2 & 0.2
 \\
Spain & 1.1 & 0.5 & 0.6 & 0.5 & -0.1 & -0.2 & 0.1
 \\
Finland & 3.9 & 2.3 & 1.6 & 0.4 & -3.5 & -3.2 & -0.3
 \\
France & 10.1 & 4.3 & 5.8 & 0.4 & 1.0 & 1.3 & -0.2
 \\
Croatia & 3.4 & 1.8 & 1.6 & 0.9 & 3.6 & 2.7 & 0.9
 \\
Ireland & 11.3 & 6.7 & 4.6 & 0.3 & 4.9 & 4.9 & 0.0
 \\
Italy & 10.8 & 5.0 & 5.8 & 0.6 & -1.0 & -0.8 & -0.2
 \\
Lithuania & 10.3 & 7.5 & 2.8 & 0.6 & 11.2 & 11.6 & -0.3
 \\
Luxembourg & 6.4 & 4.4 & 2.0 & 1.1 & 1.5 & 2.1 & -0.6
 \\
Latvia & 25.4 & 11.2 & 14.3 & 2.0 & 44.1 & 26.3 & 17.9
 \\
Malta & 1.2 & 0.0 & 1.2 & 1.2 & 0.3 & 0.0 & 0.3
 \\
Netherlands & 4.7 & 2.8 & 2.0 & 0.6 & 0.1 & 0.5 & -0.4
 \\
Portugal & 3.7 & 1.2 & 2.5 & 0.5 & 3.4 & 1.8 & 1.6
 \\
Slovenia & 7.4 & 4.0 & 3.4 & 2.2 & 4.4 & 3.8 & 0.6
 \\
Slovakia & 8.6 & 4.0 & 4.6 & 0.7 & 0.1 & 0.4 & -0.3
 \\
\midrule
\textbf{Euro area} & \textbf{7.8} & \textbf{3.6} & \textbf{4.2} & \textbf{0.6} & \textbf{1.2} & \textbf{1.1} & \textbf{0.1}
 \\
\bottomrule
\end{tabular}
\par\vspace{0.35em}\footnotesize\noindent Sources: Eurostat HICP, HBS, FIGARO, authors' calculations.
\par\vspace{0.25em}\footnotesize\noindent Notes. Entries are cumulative percentage-point effects of observed energy-price increases between June 2021 and June 2023. The direct shocks are country-specific cumulative HICP changes mapped to FIGARO energy sectors: upstream oil and gas extraction (B), refined petroleum products (C19), and electricity/gas utilities (D35). The price response is computed as $\Delta p = (I - A^{\prime})^{-1}s$, with $A$ built from the latest FIGARO matrix used in the simulations. The direct effect is the component from the shocked energy sectors in the household consumption basket. Indirect is the Leontief propagation component. Average inflation total includes direct and indirect effects before rounding. Income inflation gap = Q1 - Q5.
\end{table}
\vspace{0.6em}
\begin{table}[H]
\centering
\caption{Effect of a 10 Percent Shock to Cereal Prices, 2023}
\scriptsize
\setlength{\tabcolsep}{1.2pt}
\begin{tabular}{lrrrrrrr}
\toprule
Country & \shortstack{Effect on average\\inflation -- Total} & \shortstack{Direct effect on\\average inflation} & \shortstack{Indirect effect on\\average inflation} & \shortstack{Indirect effect on average\\inflation through trade} & \shortstack{Effect on\\inflation gap} & \shortstack{Direct effect on\\inflation gap} & \shortstack{Indirect effect on\\inflation gap}
 \\
\midrule
Austria & 0.7 & 0.2 & 0.5 & 0.1 & -0.0 & 0.0 & -0.0
 \\
Belgium & 1.0 & 0.3 & 0.7 & 0.3 & -0.1 & 0.0 & -0.1
 \\
Cyprus & 1.3 & 0.4 & 0.8 & 0.0 & -0.1 & 0.0 & -0.1
 \\
Germany & 0.7 & 0.2 & 0.4 & 0.1 & -0.0 & 0.0 & -0.0
 \\
Estonia & 1.4 & 0.4 & 0.9 & 0.1 & -0.1 & -0.0 & -0.1
 \\
Greece & 1.3 & 0.5 & 0.8 & 0.0 & -0.1 & -0.0 & -0.1
 \\
Spain & 1.0 & 0.3 & 0.7 & 0.1 & -0.0 & 0.0 & -0.0
 \\
Finland & 1.2 & 0.3 & 0.9 & 0.1 & -0.1 & 0.0 & -0.1
 \\
France & 1.1 & 0.3 & 0.8 & 0.1 & -0.1 & 0.0 & -0.1
 \\
Croatia & 1.1 & 0.4 & 0.7 & 0.1 & -0.1 & -0.0 & -0.1
 \\
Ireland & 0.8 & 0.3 & 0.5 & 0.0 & -0.0 & -0.0 & -0.0
 \\
Italy & 1.1 & 0.5 & 0.6 & 0.1 & 0.2 & 0.1 & 0.1
 \\
Lithuania & 1.3 & 0.6 & 0.7 & 0.1 & -0.1 & 0.0 & -0.1
 \\
Luxembourg & 0.5 & 0.2 & 0.3 & 0.1 & -0.0 & -0.0 & -0.0
 \\
Latvia & 1.6 & 0.6 & 1.0 & 0.2 & -0.1 & -0.0 & -0.1
 \\
Malta & 0.6 & 0.4 & 0.2 & 0.1 & -0.0 & -0.0 & -0.0
 \\
Netherlands & 0.9 & 0.3 & 0.6 & 0.1 & -0.1 & 0.0 & -0.1
 \\
Portugal & 1.2 & 0.3 & 0.9 & 0.2 & -0.1 & 0.0 & -0.1
 \\
Slovenia & 0.9 & 0.5 & 0.4 & 0.1 & -0.0 & 0.0 & -0.0
 \\
Slovakia & 1.1 & 0.5 & 0.6 & 0.1 & -0.3 & -0.1 & -0.1
 \\
\midrule
\textbf{Euro area} & \textbf{1.0} & \textbf{0.3} & \textbf{0.6} & \textbf{0.1} & \textbf{-0.0} & \textbf{0.0} & \textbf{-0.0}
 \\
\bottomrule
\end{tabular}
\par\vspace{0.35em}\footnotesize\noindent Sources: Eurostat HICP, HBS, FIGARO, authors' calculations.
\par\vspace{0.25em}\footnotesize\noindent Notes. Entries are percentage-point effects of the sectoral input price shock shown in the table title. The price response is computed as $\Delta p = (I - A^{\prime})^{-1}s$, with $A$ built from FIGARO. The direct effect is the component from the shocked sector in the household consumption basket. Indirect is the Leontief propagation component. Average inflation total includes direct and indirect effects before rounding. Income inflation gap = Q1 - Q5 of revenues.
\end{table}
\vspace{0.6em}
\begin{table}[H]
\centering
\caption{Effect of a 40\% Oil and Gas Shock with Cross-Country Propagation, June 2021 June 2023}
\scriptsize
\setlength{\tabcolsep}{1.2pt}
\begin{tabular}{lrrrrrrr}
\toprule
Country & \shortstack{Effect on average\\inflation -- Total} & \shortstack{Direct effect on\\average inflation} & \shortstack{Indirect effect on\\average inflation} & \shortstack{Indirect effect on average\\inflation through trade} & \shortstack{Effect on\\inflation gap} & \shortstack{Direct effect on\\inflation gap} & \shortstack{Indirect effect on\\inflation gap}
 \\
\midrule
Austria & 11.8 & 3.9 & 7.9 & 1.4 & 4.5 & 2.0 & 2.5
 \\
Belgium & 9.8 & 6.2 & 3.7 & 2.0 & 0.9 & 1.5 & -0.6
 \\
Cyprus & 7.5 & 4.7 & 2.8 & 0.8 & 4.0 & 3.6 & 0.3
 \\
Germany & 7.4 & 4.4 & 3.0 & 0.5 & 1.1 & 1.5 & -0.3
 \\
Estonia & 12.2 & 6.6 & 5.6 & 1.2 & 11.8 & 9.3 & 2.4
 \\
Greece & 6.9 & 4.0 & 2.9 & 0.3 & 3.6 & 3.1 & 0.5
 \\
Spain & 7.4 & 4.5 & 3.0 & 0.6 & 2.3 & 2.1 & 0.2
 \\
Finland & 5.7 & 3.6 & 2.1 & 0.4 & -5.0 & -4.6 & -0.4
 \\
France & 9.3 & 4.5 & 4.7 & 0.6 & 0.7 & 1.0 & -0.3
 \\
Croatia & 11.1 & 7.0 & 4.1 & 1.2 & 14.6 & 11.7 & 3.0
 \\
Ireland & 6.0 & 4.0 & 2.0 & 0.3 & 2.4 & 2.4 & -0.0
 \\
Italy & 9.2 & 4.9 & 4.3 & 0.5 & -0.8 & -0.7 & -0.1
 \\
Lithuania & 8.4 & 6.5 & 1.9 & 0.6 & 6.2 & 6.4 & -0.2
 \\
Luxembourg & 7.7 & 6.0 & 1.8 & 1.2 & 1.0 & 1.5 & -0.5
 \\
Latvia & 14.1 & 7.2 & 6.9 & 1.8 & 21.1 & 13.0 & 8.1
 \\
Malta & 6.1 & 2.9 & 3.2 & 1.1 & 1.7 & 1.2 & 0.5
 \\
Netherlands & 6.4 & 4.3 & 2.0 & 0.7 & 0.1 & 0.6 & -0.4
 \\
Portugal & 8.8 & 3.7 & 5.2 & 1.4 & 6.7 & 3.5 & 3.2
 \\
Slovenia & 8.3 & 5.8 & 2.5 & 1.4 & 4.6 & 4.2 & 0.4
 \\
Slovakia & 12.5 & 6.3 & 6.2 & 0.6 & 0.1 & 0.6 & -0.4
 \\
\midrule
\textbf{Euro area} & \textbf{8.4} & \textbf{4.6} & \textbf{3.8} & \textbf{0.7} & \textbf{1.3} & \textbf{1.2} & \textbf{0.0}
 \\
\bottomrule
\end{tabular}
\par\vspace{0.35em}\footnotesize\noindent Sources: Eurostat HICP, HBS, FIGARO, authors' calculations.
\par\vspace{0.25em}\footnotesize\noindent Notes. Entries are percentage-point effects over June 2021--June 2023 of a uniform 40 percent shock across the oil and gas energy chain. The directly shocked FIGARO sectors are B (mining and quarrying, including NACE B06 crude petroleum and natural gas extraction), C19 (coke and refined petroleum products), and D35 (electricity, gas, steam and air-conditioning supply). The shock vector is $s_{c,k}=0.40$ for $k\in\{B,C19,D35\}$ and $s_{c,k}=0$ otherwise, for every country $c$. The price response is $\Delta p = (I - A^{\prime})^{-1}s$, where $A_{(o,i),(d,j)}=Z_{(o,i),(d,j)}/x_{d,j}$ retains bilateral intermediate flows from origin country $o$ and sector $i$ to destination country $d$ and sector $j$. Direct effects arise from household consumption weights mapped to B, C19 and D35; indirect effects include domestic and cross-country propagation through intermediate imports and exports. The trade column is the difference between the full EA20 solution and a counterfactual solution in which all off-diagonal country blocks of $A$ are set to zero; it is included in, rather than added to, the total indirect effect. Imports from origins outside the EA20 system enter production costs but are treated as exogenous and do not generate additional endogenous feedback. FIGARO flows are measured in current million euros at basic prices; the dataset used here does not provide separate import or export deflators. Income inflation gap = Q1 - Q5 of revenues.
\end{table}
\vspace{0.6em}
